5. Results
5.1 Point Forecast Accuracy
Table 3 reports MAE and RMSE for all five point forecast models and the PJM Day-Ahead external benchmark across the full-year 2014 test set and the January 6–8 vortex window. All weather-informed model results use same-hour ERA5 retrospective reanalysis weather input and reflect an idealized retrospective weather-information setting rather than an operational forecast-weather setting.
Across the full-year 2014 test set, GBoost achieves the lowest MAE (721 MW) and lowest RMSE (994 MW) among all evaluated point forecasts. Linear Regression reduces full-year MAE to 2,354 MW relative to Persistence-1h (2,745 MW), indicating that the weather and calendar features carry predictive value beyond one-hour load persistence. Naive-24h and Naive-168h produce substantially higher full-year errors (5,960 MW and 8,009 MW MAE respectively), confirming that models using only load-memory rules without weather or calendar inputs are inadequate as standalone forecasts over a full annual cycle.
PJM Day-Ahead produces a full-year MAE of 1,937 MW. GBoost produces a full-year MAE of 721 MW. This difference reflects the retrospective informational advantage conferred by same-hour ERA5 reanalysis weather inputs; it does not constitute evidence of operational superiority, as PJM Day-Ahead is produced in real time using forecast-weather inputs not controlled by this study.
During the January 6–8 vortex window, GBoost again produces the lowest MAE (1,705 MW) and RMSE (2,037 MW) among evaluated point forecasts. However, all models show error degradation relative to their full-year performance, indicating that the cold-event window presents greater forecasting difficulty across the board. The degradation is most severe for the seasonal naive baselines: Naive-24h MAE increases from 5,960 MW full-year to 15,598 MW during the vortex window, and Naive-168h MAE increases from 8,009 MW to 24,531 MW. Linear Regression vortex MAE (4,248 MW) is approximately 1.8 times its full-year MAE (2,354 MW). GBoost vortex MAE (1,705 MW) is approximately 2.4 times its full-year MAE (721 MW). PJM Day-Ahead vortex MAE (3,148 MW) is approximately 1.6 times its full-year MAE (1,937 MW). The relative degradation pattern across models is discussed further in Section 6.
5.2 Event-Peak Point Forecast Errors
The January 6–8, 2014 Polar Vortex produced a near-annual-peak RTO load of 140,510.2 MW at Jan 7 18:00 EPT, equal to 99.18% of the 2014 annual peak of 141,677.9 MW recorded on Jun 17 17:00 EPT. Table 5 reports the point forecast and prediction interval results at this single peak hour.
Using the error convention $e_t = y_t - \hat{y}_t$ (positive = underprediction, negative = overprediction), the point forecast errors at Jan 7 18:00 EPT are:


Persistence-1h: predicted 137,604 MW, error +2,906 MW (underprediction)


Naive-24h: predicted 130,537 MW, error +9,973 MW (underprediction)


Naive-168h: predicted 109,742 MW, error +30,769 MW (underprediction)


Linear Regression: predicted 134,985 MW, error +5,525 MW (underprediction)


GBoost: predicted 141,898 MW, error −1,388 MW (overprediction)


PJM Day-Ahead: predicted 137,965 MW, error +2,545 MW (underprediction)


All models except GBoost underpredict the event peak. The seasonal naive baselines produce large positive errors consistent with their full-year vortex-window degradation: Naive-168h, anchored to load from the prior week rather than to the contemporaneous cold-event conditions, misses the event peak by 30,769 MW. Linear Regression underpredicts by 5,525 MW despite using the same ERA5 retrospective weather inputs as GBoost, suggesting that the linear specification was less able than GBoost to represent event-window behavior under extreme cold conditions. GBoost slightly overpredicts the event peak by 1,388 MW. PJM Day-Ahead underpredicts by 2,545 MW.
These single-hour results are not sufficient to characterize overall model behavior and should be read alongside the window-level MAE and RMSE in Table 3.
5.3 Probabilistic Forecast Performance
Table 4 reports QR-GBT probabilistic forecast metrics across the full-year 2014 test set, the January 6–8 vortex window, the June 16–18 summer near-peak window, and the top 1% load-hour subset. All probabilistic evaluation uses post-hoc rearranged quantile estimates as described in Section 3.5; reported prediction intervals should be interpreted in light of the high quantile crossing rate (5,786 of 8,760 full-year test hours, 66%) prior to rearrangement.
Full-year 2014. The QR-GBT q50 median forecast achieves a full-year MAE of 761 MW, close to the GBoost point forecast MAE of 721 MW. The nominal 90% prediction interval (q05–q95) achieves an empirical coverage rate of 86.8% against a nominal target of 90%. The nominal 98% prediction interval (q01–q99) achieves an empirical coverage rate of 97.3%. The mean pinball loss averaged across the seven estimated quantile levels is 153, and the mean Winkler score for the nominal 90% interval is 4,627 MW.
January 6–8 vortex window. During the vortex window, the q50 MAE increases to 2,130 MW. The nominal 90% PI empirical coverage falls to 66.7%, 20.1 percentage points below its full-year rate and 23.3 percentage points below the nominal 90% target. The nominal 98% PI empirical coverage falls to 84.7%, 12.6 percentage points below its full-year rate and 13.3 percentage points below the nominal 98% target. The mean pinball loss increases to 473 and the mean Winkler score for the nominal 90% interval increases to 15,476 MW, indicating poorer joint interval sharpness and coverage performance relative to the full-year evaluation.
June 16–18 summer near-peak window. The q50 MAE during the summer near-peak window is 1,060 MW. The nominal 90% PI coverage is 86.1% and the nominal 98% PI coverage is 93.1%. The mean pinball loss is 209 and the Winkler score for the nominal 90% interval is 6,355 MW.
Top 1% load hours. The top 1% of 2014 test hours by observed load (approximately 88 hours) show 90% PI coverage of 63.2% and 98% PI coverage of 82.8%. The q50 MAE for this subset is 2,539 MW, the mean pinball loss is 502, and the Winkler score for the nominal 90% interval is 16,168 MW. These results are consistent with the vortex-window findings and indicate that coverage degradation is not limited to the polar vortex event also appears in the top 1% load-hour subset of the 2014 evaluation year.
Figure 3 shows the QR-GBT retrospective probabilistic forecast during the January 6–8 vortex window, including the nominal 90% and 98% prediction interval bands and the observed load series. As shown in Figure 3, the event peak at Jan 7 18:00 EPT falls outside the nominal 98% prediction interval. All forecast values in Figure 3 use same-hour ERA5 retrospective reanalysis weather input.
5.4 Calibration Breakdown Across Evaluation Windows
Figure 4 summarizes empirical coverage rates across evaluation windows for the nominal 90% and 98% prediction intervals. The full-year 90% PI coverage rate of 86.8% is reasonably close to the nominal 90% target, which might suggest adequate probabilistic calibration under a single-number summary. However, disaggregation by evaluation window reveals that this aggregate rate masks substantial undercoverage during periods of high and extreme load.
The nominal 90% PI coverage falls from 86.8% (full-year) to 66.7% during the vortex window — a decline of 20.1 percentage points. The nominal 98% PI coverage falls from 97.3% (full-year) to 84.7% during the vortex window. The top 1% load-hour subset shows 90% PI coverage of 63.2% and 98% PI coverage of 82.8%, similar in magnitude to the vortex window findings. The summer near-peak window 90% PI coverage (86.1%) is close to the full-year rate and does not show the same pattern of degradation.
These results indicate that aggregate full-year calibration metrics do not adequately characterize probabilistic forecast performance during the high-load and extreme cold-weather conditions that are operationally most consequential. The mechanisms underlying this calibration breakdown are examined in Section 6.
5.5 Winter versus Summer Near-Peak Comparison
Figure 5 places the January 6–8 vortex window results alongside the June 16–18 summer near-peak window results. The two windows have similar peak-load magnitudes: the winter vortex peak is 140,510 MW and the summer peak is 141,678 MW, a difference of 1,168 MW or less than 1% of the summer peak. Despite this similarity in peak magnitude, probabilistic performance is materially worse during the cold-event window than during the summer near-peak window across all reported metrics.
The q50 MAE during the vortex window (2,130 MW) is approximately twice the q50 MAE during the summer window (1,060 MW). The mean pinball loss during the vortex window (473) is approximately twice that of the summer window (209). The nominal 90% PI Winkler score during the vortex window (15,476 MW) is approximately 2.4 times that of the summer window (6,355 MW). Nominal 90% PI coverage is 66.7% during the vortex window and 86.1% during the summer window; nominal 98% PI coverage is 84.7% and 93.1% respectively.
The comparison indicates that peak-load magnitude alone does not account for the degradation in probabilistic performance during the vortex window. The additional difficulty associated with the cold-event window, relative to a summer near-peak window of comparable load magnitude, is addressed in the Discussion.
5.6 Tail-Risk Event-Peak Interval Check
Table 5 and Figure 6 report the QR-GBT prediction interval bounds at the event peak hour (Jan 7 18:00 EPT, observed load 140,510 MW). All reported values use post-hoc rearranged quantile estimates as described in Section 3.5.
At the event peak, the QR-GBT q50 estimate is 135,357 MW, reflecting a median forecast underprediction of 5,153 MW. The post-rearrangement q95 estimate is 139,410 MW and the post-rearrangement q99 estimate is 139,585 MW. At the event peak, the observed load exceeded the post-rearrangement QR-GBT q99 estimate of 139,585 MW, leaving the event peak outside the nominal 98% prediction interval. The q99 estimate falls 925 MW below the observed peak load of 140,510 MW.
The event peak also falls above the post-rearrangement q95 bound of 139,410 MW, confirming that the observed load at this hour lies outside the nominal 90% prediction interval as well. Figure 6 illustrates the relationship between the observed event peak and the estimated quantile bounds at this hour.
These results indicate that at the hour of maximum cold-event load, the QR-GBT upper quantile estimates were insufficient to bound the observed outcome even at the 99th percentile level. The interpretation of this finding in terms of training-distribution coverage and the retrospective weather-input setting is addressed in Section 6.